\begin{abcliste}
\abc
Vor P zieht die Schnur den Stein zum Pfosten hin. Diese Kraft zeigt immer quer zur Bewegung und krümmt die Bahn ständig. Nach P zieht die Schnur nicht mehr, und keine andere horizontale Kraft wirkt. Nach dem ersten Newtonschen Gesetz bewegt sich der Stein deshalb geradlinig weiter, in der Richtung, die er in P hatte, also entlang der Tangente:
\[
 \result{\text{Bahn B}}
\]
Die anderen Bahnen bräuchten eine Kraft:
\begin{itemize}
\item C (weiter im Bogen, langsam nach aussen): Um weiter zur Kreismitte hin abzubiegen, bräuchte es eine Kraft zum Zentrum. Der Stein „erinnert“ sich nicht an den Kreis.
\item D (geradlinig nach aussen) und A (nach aussen gekrümmt): Dafür bräuchte es eine Kraft nach aussen. Eine solche „Zentrifugalkraft“ gibt es nicht: Die Schnur hat nur nach innen gezogen.
\end{itemize}

\abc
\result{\text{Keine.}} Die Gewichtskraft und die Normalkraft des Eises sind senkrecht und heben sich auf; die Reibung ist vernachlässigbar. Es gibt weder eine Kraft in Bewegungsrichtung noch eine Kraft nach aussen, und auch keine Kraft zum Zentrum, die langsam abklingt: Die Schnur kann nur ziehen, solange sie den Stein festhält.

\abc
Die Schnur liefert die Kraft zum Zentrum hin, die die Richtung der Geschwindigkeit ständig ändert; ohne sie würde der Stein schon vorher geradeaus gleiten. Nach P wirkt keine horizontale Kraft mehr, also \result{\text{bleibt die Geschwindigkeit nach Betrag und Richtung gleich}}: Der Stein gleitet mit der Geschwindigkeit, die er in P hatte, geradlinig weiter.
\end{abcliste}
